% ============================================================
\chapter{Th\'eorie des Points Fixes en Topologie Avanc\'ee}
\label{ch:points_fixes}
% ============================================================

Les th\'eor\`emes de points fixes prennent une saveur particuli\`ere dans le cadre de la topologie non-Hausdorff. Ici, pas de m\'etrique, pas de contraction~: les points fixes \'emergent de la structure d'ordre et de la compl\'etude dirig\'ee. Le th\'eor\`eme de Tarski (1955) affirme que toute application monotone sur un treillis complet poss\`ede un point fixe~; celui de Kleene le construit par it\'eration d\'enombrable dans les $\omega$-cpo~; et le th\'eor\`eme de Pataraia (1997), longtemps rest\'e confidentiel dans la litt\'erature g\'eorgienne, \'etend ces r\'esultats aux dcpo g\'en\'eraux. Ces th\'eor\`emes sont les fondations de la s\'emantique d\'enotationnelle~: d\'efinir la signification d'un programme r\'ecursif, c'est pr\'ecis\'ement trouver un point fixe dans un domaine de Scott.

\begin{intuition}
Les th\'eor\`emes de points fixes jouent un r\^ole central en informatique th\'eorique et
en th\'eorie des domaines. Contrairement aux th\'eor\`emes classiques de Banach ou Brouwer,
qui reposent sur des structures m\'etriques ou topologiques classiques, les th\'eor\`emes
pr\'esent\'es ici exploitent l'ordre et la compl\'etude dirig\'ee. Le th\'eor\`eme de
Tarski garantit l'existence de points fixes dans les treillis complets, celui de Kleene
les construit effectivement par it\'eration, et le th\'eor\`eme de Pataraia \'etend ces
r\'esultats aux dcpos sans hypoth\`ese de d\'enombrabilit\'e.
\end{intuition}

% ------------------------------------------------------------
\section{Th\'eor\`eme de Tarski pour les treillis complets}
% ------------------------------------------------------------

\begin{definition}[Treillis complet]
\index{treillis complet}
Un ensemble ordonn\'e $(L, \leq)$ est un \emph{treillis complet} si toute partie
$S \subseteq L$ admet un supremum $\bigvee S$ et un infimum $\bigwedge S$. En
particulier, $L$ poss\`ede un plus petit \'el\'ement $\bot = \bigvee \emptyset$ et un
plus grand \'el\'ement $\top = \bigwedge \emptyset$.
\end{definition}

\begin{theoreme}[Knaster--Tarski]
\label{thm:knaster-tarski}
\index{th\'eor\`eme de Knaster--Tarski}
Soit $(L, \leq)$ un treillis complet et $f\colon L \to L$ une fonction croissante.
Alors~:
\begin{enumerate}
  \item $f$ poss\`ede un plus petit point fixe
    $\mathrm{lfp}(f) = \bigwedge \{x \in L : f(x) \leq x\}$.
  \item $f$ poss\`ede un plus grand point fixe
    $\mathrm{gfp}(f) = \bigvee \{x \in L : x \leq f(x)\}$.
  \item L'ensemble des points fixes $\mathrm{Fix}(f)$ est un treillis complet.
\end{enumerate}
\end{theoreme}

\begin{proof}
(1) Posons $P = \{x \in L : f(x) \leq x\}$ et $a = \bigwedge P$. Pour tout $x \in P$,
$a \leq x$, donc $f(a) \leq f(x) \leq x$ par croissance. Ainsi $f(a)$ est un minorant
de $P$, d'o\`u $f(a) \leq a$. Ceci montre que $a \in P$.

Par croissance, $f(f(a)) \leq f(a)$, donc $f(a) \in P$, d'o\`u $a \leq f(a)$. Combin\'e
avec $f(a) \leq a$, on obtient $f(a) = a$.

Si $b$ est un autre point fixe, $b \in P$ donc $a \leq b$~: $a$ est bien le plus petit
point fixe.

(2) Dual de (1).

(3) Soit $S \subseteq \mathrm{Fix}(f)$. L'ensemble $\{x \in L : x \geq s \; \forall s \in S,
\; f(x) \leq x\}$ est non vide (il contient $\top$). Son infimum est un point fixe de $f$
qui est le supremum de $S$ dans $\mathrm{Fix}(f)$.
\end{proof}

\begin{remarque}
Le th\'eor\`eme de Knaster--Tarski ne requiert aucune hypoth\`ese de continuit\'e~: la
seule monotonie suffit. En revanche, il ne fournit pas de m\'ethode constructive pour
calculer le point fixe.
\end{remarque}

\begin{exemple}[Application aux syst\`emes d'\'equations]
Soit $\Sigma$ un alphabet et $\mathcal{P}(\Sigma^*)$ le treillis des langages sur $\Sigma$.
Un syst\`eme d'\'equations de langages $X_i = f_i(X_1, \ldots, X_n)$ o\`u les $f_i$ sont
croissantes admet une plus petite solution, par le th\'eor\`eme de Tarski appliqu\'e au
treillis $\mathcal{P}(\Sigma^*)^n$.
\end{exemple}

% ------------------------------------------------------------
\section{Th\'eor\`eme de Kleene}
% ------------------------------------------------------------

\begin{definition}[dcpo point\'e]
\index{dcpo point\'e}
Un \emph{dcpo point\'e} est un dcpo $(D, \leq)$ poss\'edant un plus petit \'el\'ement $\bot$.
\end{definition}

\begin{theoreme}[Kleene]
\label{thm:kleene}
\index{th\'eor\`eme de Kleene}
Soit $(D, \leq)$ un dcpo point\'e et $f\colon D \to D$ une fonction Scott-continue. Alors
$f$ poss\`ede un plus petit point fixe, et celui-ci est donn\'e par~:
\[
  \mathrm{lfp}(f) = \bigvee_{n \geq 0} f^n(\bot).
\]
\end{theoreme}

\begin{proof}
Montrons d'abord que la suite $(f^n(\bot))_{n \geq 0}$ est croissante. On a
$\bot \leq f(\bot)$ car $\bot$ est le plus petit \'el\'ement. Par r\'ecurrence, si
$f^n(\bot) \leq f^{n+1}(\bot)$, alors par croissance de $f$,
$f^{n+1}(\bot) \leq f^{n+2}(\bot)$.

L'ensemble $\{f^n(\bot) : n \geq 0\}$ est dirig\'e (c'est une cha\^ine), donc son
supremum $a = \bigvee_n f^n(\bot)$ existe dans $D$.

Par Scott-continuit\'e~:
\[
  f(a) = f\!\left(\bigvee_{n \geq 0} f^n(\bot)\right) = \bigvee_{n \geq 0} f^{n+1}(\bot) = a.
\]
Donc $a$ est un point fixe.

Si $b$ est un autre point fixe, alors $\bot \leq b$ et par r\'ecurrence $f^n(\bot) \leq b$
pour tout $n$, d'o\`u $a = \bigvee_n f^n(\bot) \leq b$.
\end{proof}

\begin{attention}
Le th\'eor\`eme de Kleene n\'ecessite la Scott-continuit\'e (pas seulement la
croissance). Une fonction croissante sur un dcpo point\'e peut ne pas avoir de plus petit
point fixe obtenu par it\'eration depuis $\bot$.
\end{attention}

\begin{corollaire}[It\'eration transfinie]
Si $f\colon L \to L$ est une fonction croissante sur un treillis complet, le plus petit
point fixe est atteint par it\'eration transfinie~:
\[
  f^\alpha(\bot) = \begin{cases}
    f(f^{\alpha-1}(\bot)) & \text{si } \alpha \text{ est successeur}, \\
    \bigvee_{\beta < \alpha} f^\beta(\bot) & \text{si } \alpha \text{ est limite}.
  \end{cases}
\]
Il existe un ordinal $\alpha_0$ tel que $f^{\alpha_0}(\bot) = \mathrm{lfp}(f)$.
\end{corollaire}

\begin{exemple}[S\'emantique des boucles]
La s\'emantique d'une boucle \texttt{while b do c} est le plus petit point fixe de
l'op\'erateur $\Phi\colon [S \to S_\bot] \to [S \to S_\bot]$ d\'efini par~:
\[
  \Phi(g)(s) = \begin{cases}
    g(\llbracket c \rrbracket(s)) & \text{si } \llbracket b \rrbracket(s) = \mathtt{true}, \\
    s & \text{si } \llbracket b \rrbracket(s) = \mathtt{false}, \\
    \bot & \text{sinon}.
  \end{cases}
\]
Par le th\'eor\`eme de Kleene, $\mathrm{lfp}(\Phi) = \bigvee_n \Phi^n(\bot_{S \to S_\bot})$.
\end{exemple}

% ------------------------------------------------------------
\section{Th\'eor\`eme de Pataraia}
% ------------------------------------------------------------

\begin{theoreme}[Pataraia]
\label{thm:pataraia}
\index{th\'eor\`eme de Pataraia}
Soit $(D, \leq)$ un dcpo point\'e et $f\colon D \to D$ une fonction croissante (pas
n\'ecessairement continue) telle que $f(x) \geq x$ pour tout $x$ (\emph{inflationnaire}).
Alors $f$ poss\`ede un point fixe.
\end{theoreme}

\begin{proof}
D\'efinissons l'ensemble $A = \{ x \in D : \bot \leq x \leq f(x) \text{ et } x \text{ est
le supremum d'une cha\^ine d'it\'er\'es de } f \text{ depuis } \bot \}$. Plus
pr\'ecis\'ement, consid\'erons l'ensemble $I$ de tous les sous-ensembles $S$ de $D$ tels
que~:
\begin{enumerate}
  \item $\bot \in S$~;
  \item $S$ est stable par $f$~;
  \item $S$ est stable par suprema dirig\'es (dans $D$).
\end{enumerate}
Soit $C = \bigcap I$ le plus petit tel ensemble. On montre que $C$ est une cha\^ine (en
v\'erifiant que l'ensemble des \'el\'ements comparables \`a tous les \'el\'ements de $C$
satisfait les trois conditions). Le supremum $c = \bigvee C$ existe dans $D$ et v\'erifie
$f(c) \leq c$ (car $f(c) \in C$ par stabilit\'e, d'o\`u $f(c) \leq c$). Comme $f$ est
inflationnaire, $c \leq f(c)$, d'o\`u $f(c) = c$.
\end{proof}

\begin{remarque}
Le th\'eor\`eme de Pataraia g\'en\'eralise le th\'eor\`eme de Bourbaki--Witt (qui suppose
que toute cha\^ine a un supremum). L'avantage du th\'eor\`eme de Pataraia est qu'il ne
n\'ecessite que la compl\'etude dirig\'ee, pas la compl\'etude par cha\^ines.
\end{remarque}

\begin{corollaire}
Tout dcpo point\'e dans lequel toute fonction croissante inflationnaire a un point fixe
est un treillis complet.
\end{corollaire}

% ------------------------------------------------------------
\section{Plus petit point fixe en th\'eorie des domaines}
% ------------------------------------------------------------

\begin{theoreme}[Existence du plus petit point fixe]
\label{thm:lfp-domain}
Soit $D$ un domaine continu point\'e et $f\colon D \to D$ une fonction Scott-continue.
Alors~:
\begin{enumerate}
  \item Le plus petit point fixe $\mathrm{lfp}(f)$ existe et est donn\'e par la formule
    de Kleene $\mathrm{lfp}(f) = \bigvee_n f^n(\bot)$.
  \item Si de plus $D$ est un domaine alg\'ebrique effectivement pr\'esent\'e et $f$ est
    effectivement donn\'ee, alors $\mathrm{lfp}(f)$ est calculable.
  \item L'op\'erateur $\mathrm{lfp}\colon [D \to D] \to D$ est lui-m\^eme Scott-continu.
\end{enumerate}
\end{theoreme}

\begin{proof}
(1) C'est le th\'eor\`eme de Kleene (Th\'eor\`eme~\ref{thm:kleene}).

(2) Les approximations finies $f^n(\bot)$ sont calculables, et les \'el\'ements compacts
en dessous de $\mathrm{lfp}(f)$ sont \'enum\'erables.

(3) Soit $(f_i)_{i \in I}$ une famille dirig\'ee dans $[D \to D]$ et $g = \bigvee_i f_i$.
Alors~:
\[
  \mathrm{lfp}(g) = \bigvee_n g^n(\bot) = \bigvee_n \left(\bigvee_i f_i\right)^n(\bot)
  = \bigvee_i \bigvee_n f_i^n(\bot) = \bigvee_i \mathrm{lfp}(f_i).
\]
L'\'echange des suprema est justifi\'e par la Scott-continuit\'e de la composition et de
l'application.
\end{proof}

% ------------------------------------------------------------
\section{Applications aux d\'efinitions r\'ecursives}
% ------------------------------------------------------------

\begin{proposition}[D\'efinitions r\'ecursives de fonctions]
Soit $D$ un dcpo point\'e. Une d\'efinition r\'ecursive de la forme
\[
  f(x) = \Phi(f, x)
\]
o\`u $\Phi\colon [D \to D] \times D \to D$ est Scott-continue, poss\`ede une plus petite
solution $f^* \in [D \to D]$, donn\'ee par $f^* = \mathrm{lfp}(\lambda g.\, \lambda x.\, \Phi(g,x))$.
\end{proposition}

\begin{exemple}[Factorielle]
La d\'efinition r\'ecursive
\[
  \mathtt{fact}(n) = \begin{cases} 1 & \text{si } n = 0, \\ n \cdot \mathtt{fact}(n-1) & \text{sinon} \end{cases}
\]
correspond \`a $\mathtt{fact} = \mathrm{lfp}(\Phi)$ o\`u $\Phi(g)(n) = \text{si } n = 0
\text{ alors } 1 \text{ sinon } n \cdot g(n-1)$. L'op\'erateur $\Phi$ est Scott-continu
sur $[\N_\bot \to \N_\bot]$.
\end{exemple}

\begin{theoreme}[\'Equations de domaines r\'ecursives]
\index{\'equations de domaines}
Si $F$ est un foncteur continu localement continu sur la cat\'egorie des domaines continus
avec paires de r\'etraction, alors l'\'equation de domaine
\[
  D \cong F(D)
\]
admet une solution, obtenue comme limite projective (et aussi colimite inductive) de la
suite $D_0 \leftarrow D_1 \leftarrow \cdots$ o\`u $D_{n+1} = F(D_n)$.
\end{theoreme}

\begin{formules}[title=\textbf{R\'ecapitulatif des th\'eor\`emes de points fixes}]
\begin{center}
\renewcommand{\arraystretch}{1.3}
\begin{tabular}{|l|l|l|l|}
\hline
\textbf{Th\'eor\`eme} & \textbf{Structure} & \textbf{Hypoth\`ese sur $f$} & \textbf{R\'esultat} \\
\hline
Knaster--Tarski & Treillis complet & Croissante & $\mathrm{Fix}(f)$ treillis complet \\
Kleene & dcpo point\'e & Scott-continue & $\mathrm{lfp}(f) = \bigvee f^n(\bot)$ \\
Pataraia & dcpo point\'e & Croissante inflat. & Existence de point fixe \\
Bourbaki--Witt & Ordonn\'e complet & Inflationnaire & Existence de point fixe \\
\hline
\end{tabular}
\end{center}
\end{formules}

% ------------------------------------------------------------
\section{Exercices}
% ------------------------------------------------------------

\begin{exercice}
Soit $L = \mathcal{P}(\N)$ le treillis des parties de $\N$. D\'eterminer le plus petit
et le plus grand point fixe de $f(A) = \N \setminus A$.
\end{exercice}

\begin{exercice}
Montrer que dans le th\'eor\`eme de Kleene, la suite $(f^n(\bot))$ peut se stabiliser apr\`es
un nombre fini d'\'etapes si et seulement si $f$ est \og{}finie\fg au sens o\`u le plus
petit point fixe est un \'el\'ement compact du dcpo.
\end{exercice}

\begin{exercice}
Donner un exemple de dcpo point\'e $D$ et de fonction croissante $f\colon D \to D$ telle
que la suite $(f^n(\bot))$ ne converge pas vers le plus petit point fixe de $f$ (c'est-\`a-dire
telle que $\bigvee_n f^n(\bot) < \mathrm{lfp}(f)$).
\end{exercice}

\begin{exercice}
Soit $D$ un dcpo point\'e et $f, g\colon D \to D$ deux fonctions Scott-continues telles que
$f \leq g$ point par point. Montrer que $\mathrm{lfp}(f) \leq \mathrm{lfp}(g)$.
\end{exercice}

\begin{exercice}[Th\'eor\`eme de Bekic]
Soient $D, E$ deux dcpos point\'es et $f\colon D \times E \to D$, $g\colon D \times E \to E$
des fonctions Scott-continues. Montrer que le plus petit point fixe de $(f,g)\colon D \times E
\to D \times E$ est donn\'e par $(\mathrm{lfp}(\hat{f}), \mathrm{lfp}(\hat{g}))$ o\`u
$\hat{f}$ et $\hat{g}$ sont d\'efinies par substitution mutuelle.
\end{exercice}

\begin{exercice}
Soit $\Sigma = \{a, b\}$ et $\mathcal{P}(\Sigma^*)$ le treillis des langages.
L'op\'erateur $\Phi(X) = \{\varepsilon\} \cup \{a\} \cdot X \cdot \{b\}$ est croissant.
Calculer $\mathrm{lfp}(\Phi)$ par le th\'eor\`eme de Kleene.
\end{exercice}

\begin{exercice}
Montrer que le th\'eor\`eme de Pataraia implique le th\'eor\`eme de Zorn~: tout ensemble
ordonn\'e inductif (toute cha\^ine est born\'ee) admet un \'el\'ement maximal.
\emph{Indication~: consid\'erer un dcpo point\'e et une fonction de choix.}
\end{exercice}
